
\subsection{Arquitectura del Módulo de Progreso y Recomendaciones}

El módulo de progreso y recomendaciones centraliza el seguimiento
longitudinal del desempeño del alumno y expone al tutor las herramientas
de supervisión grupal necesarias para tomar decisiones pedagógicas
informadas. A diferencia de los módulos anteriores, cuyo alcance es
transaccional —un intento, un ejercicio, una autenticación—, este módulo
opera sobre el historial acumulado del alumno, sintetizando múltiples
análisis en indicadores agregados que reflejan su evolución en el tiempo.

La arquitectura del módulo se divide en tres componentes principales:

\begin{enumerate}
    \item \textbf{Capa de Presentación (Front-End):} Compuesta por el
    componente \texttt{ProgresoAlumnos.jsx} en el panel del tutor, que
    presenta gráficas de dona y de barras con la distribución del grupo
    por categoría de desempeño, y por el componente
    \texttt{PerfilAlumno.jsx} en el panel del alumno, que expone sus
    datos personales y permite actualizarlos. La clasificación del grupo
    se presenta en dos dimensiones independientes —ortografía y
    legibilidad— seleccionables mediante pestañas, con soporte
    responsivo para pantallas móviles.

    \item \textbf{Capa de Lógica de Negocio (Back-End):} Implementada
    en los controladores \texttt{alumnos.py} y \texttt{perfil.py}
    mediante FastAPI. El primero gestiona el registro, la consulta y la
    baja de alumnos, así como la clasificación del grupo por nivel de
    desempeño y la generación de datos para las gráficas de progreso.
    El segundo gestiona las operaciones de actualización de perfil tanto
    del tutor como del alumno.

    \item \textbf{Capa de Persistencia (Base de Datos):} El módulo opera
    principalmente sobre las tablas \texttt{Progreso\_Alumno} e
    \texttt{Historial\_Alumno}. La primera almacena el promedio acumulado
    vigente del alumno en cada dimensión de evaluación; la segunda
    registra una fotografía histórica del progreso cada vez que dicho
    promedio es actualizado, permitiendo trazar la evolución del alumno
    a lo largo del tiempo. Las tablas \texttt{Recomendacion} y
    \texttt{Actividad\_Sugerida} soportan la generación de actividades
    personalizadas a partir de los errores detectados.
\end{enumerate}

% ─────────────────────────────────────────────────────────
\subsubsection{Diseño del Modelo de Progreso}

El progreso del alumno se representa mediante dos dimensiones
independientes pero complementarias:

\begin{itemize}
    \item \textbf{Promedio ortográfico} (\texttt{promedio\_ortografia}):
    promedio acumulado del campo \texttt{ortografia\_score} de todos los
    análisis ortográficos registrados para el alumno.

    \item \textbf{Promedio de legibilidad} (calculado): promedio aritmético
    de los cuatro scores caligráficos —\texttt{alineacion\_score},
    \texttt{tamano\_letra\_score}, \texttt{espaciado\_score} e
    \texttt{inclinacion\_score}— derivados de los análisis caligráficos
    del alumno.
\end{itemize}

El \textbf{promedio general} se obtiene como la media de ambas
dimensiones y actúa como indicador sintético para la clasificación del
alumno dentro de su grupo:

\begin{equation}
    Promedio\_general = \frac{Promedio\_ortográfico +
    Promedio\_legibilidad}{2}
    \label{eq:promedio-general}
\end{equation}

Con base en el promedio general, el sistema clasifica a cada alumno en
una de tres categorías, conforme a la Tabla~\ref{tab:categorias-progreso}:

\begin{table}[H]
\centering
\caption{Categorías de clasificación de alumnos por promedio general}
\label{tab:categorias-progreso}
\begin{tabular}{|l|l|p{7cm}|}
\hline
\textbf{Categoría} & \textbf{Rango} & \textbf{Interpretación} \\
\hline
En riesgo   & $< 6.0$         & El alumno presenta deficiencias significativas en
                                 ortografía o legibilidad que requieren atención
                                 prioritaria del tutor. \\
Regular     & $[6.0, \ 8.0]$  & El alumno mantiene un desempeño aceptable con
                                 áreas de oportunidad identificadas. \\
Excelencia  & $> 8.0$         & El alumno muestra un desempeño consistentemente
                                 alto en ambas dimensiones evaluadas. \\
\hline
\end{tabular}
\end{table}
\newpage
\noindent\textbf{Nota de diseño:} los umbrales de clasificación (6.0 y
8.0) están alineados con la escala de evaluación del sistema educativo
mexicano de educación primaria \cite{sep2011guia}, donde 6.0 representa
el mínimo aprobatorio y 8.0 delimita el desempeño sobresaliente.

% ─────────────────────────────────────────────────────────
\subsubsection{Diseño del Flujo de Actualización de Progreso}

La tabla \texttt{Progreso\_Alumno} se actualiza cada vez que el módulo
OCR/NLP completa el análisis de un intento. El mecanismo de actualización
opera en dos pasos:

\begin{enumerate}
    \item \textbf{Actualización del promedio vigente:} al persistir los
    resultados de \texttt{Analisis\_Ortografico} y
    \texttt{Analisis\_Caligrafico}, el proceso recalcula los promedios
    acumulados del alumno mediante una consulta \texttt{UPDATE} con
    subconsulta \texttt{AVG} sobre todos sus intentos analizados hasta
    la fecha, y los escribe en \texttt{Progreso\_Alumno}.

    \item \textbf{Registro histórico:} simultáneamente, se inserta un
    nuevo registro en \texttt{Historial\_Alumno} con una instantánea
    (\textit{snapshot}) del progreso en ese momento, incluyendo la fecha
    con precisión \texttt{DATETIME2}. Este registro es inmutable y
    permite reconstruir la curva de progreso del alumno a lo largo del
    tiempo, independientemente de los cambios posteriores en
    \texttt{Progreso\_Alumno}.
\end{enumerate}

La separación entre la tabla de progreso vigente y la tabla de historial
responde a una decisión de diseño deliberada: \texttt{Progreso\_Alumno}
es la fuente de verdad para las consultas actuales del tutor y los
cálculos de clasificación, mientras que \texttt{Historial\_Alumno}
es la fuente de verdad para la visualización de la evolución temporal.
Ambas responsabilidades son distintas y justifican la existencia de dos
tablas separadas.

% ─────────────────────────────────────────────────────────
\subsubsection{Diseño del Flujo de Recomendaciones}

Cuando el análisis de un intento es completado, el sistema genera
recomendaciones de actividades personalizadas para el alumno. El flujo
es el siguiente:

\begin{enumerate}
    \item El proceso evalúa los resultados del análisis recién
    completado e identifica el tipo de error predominante: ortográfico
    si \texttt{ortografia\_score} está por debajo del umbral, caligráfico
    si alguno de los cuatro scores de legibilidad lo está.

    \item Se consulta la tabla \texttt{Actividad\_Sugerida} filtrando por
    \texttt{tipo\_error} y \texttt{tipo\_ejercicio} correspondientes al
    error identificado.

    \item Se insertan registros en la tabla \texttt{Recomendacion}
    vinculando el \texttt{id\_intento} con las actividades seleccionadas.
    Cada recomendación recibe un valor de \texttt{prioridad}
    (\texttt{alta}, \texttt{media} o \texttt{baja}) calculado con base
    en la magnitud del error: a mayor desviación respecto al umbral,
    mayor prioridad.

    \item El alumno consulta sus recomendaciones pendientes, ordenadas
    por prioridad, desde su panel. Al realizar una actividad sugerida,
    el campo \texttt{estado} de la recomendación transiciona de
    \texttt{pendiente} a \texttt{completada}.
\end{enumerate}

% ─────────────────────────────────────────────────────────
\subsubsection{Diseño de APIs}

Para soportar los flujos descritos, el diseño de la API contempla los
siguientes endpoints distribuidos entre los controladores
\texttt{alumnos.py} y \texttt{perfil.py}:

\paragraph{Endpoints de gestión de alumnos (tutor)}

\begin{table}[H]
\centering
\caption{Endpoints de gestión de alumnos}
\label{tab:endpoints-alumnos}
\begin{tabular}{|l|l|l|}
\hline
\textbf{Endpoint} & \textbf{Método} & \textbf{Acceso} \\
\hline
\texttt{/api/alumnos/registrar}              & POST & Tutor (JWT) \\
\texttt{/api/alumnos/tutor/\{id\_tutor\}}    & GET  & Tutor (JWT) \\
\texttt{/api/alumnos/baja/\{id\_alumno\}}    & PUT  & Tutor (JWT) \\
\hline
\end{tabular}
\end{table}

\paragraph{Endpoints de clasificación y progreso grupal (tutor)}

Los siguientes tres endpoints implementan la clasificación del grupo
según las categorías definidas en la Tabla~\ref{tab:categorias-progreso}.
Cada uno ejecuta la misma consulta CTE (\textit{Common Table Expression})
sobre \texttt{Progreso\_Alumno}, variando únicamente el filtro de
\texttt{promedio\_general} aplicado:

\begin{table}[H]
\centering
\caption{Endpoints de clasificación de alumnos por desempeño}
\label{tab:endpoints-clasificacion}
\begin{tabular}{|l|l|l|}
\hline
\textbf{Endpoint} & \textbf{Método} & \textbf{Acceso} \\
\hline
\texttt{/api/alumnos/riesgo/\{id\_tutor\}}      & GET & Tutor (JWT) \\
\texttt{/api/alumnos/regular/\{id\_tutor\}}     & GET & Tutor (JWT) \\
\texttt{/api/alumnos/excelencia/\{id\_tutor\}}  & GET & Tutor (JWT) \\
\hline
\end{tabular}
\end{table}

El endpoint \texttt{/api/alumnos/progreso/\{id\_tutor\}} opera mediante
el método \texttt{GET} y retorna los datos agrupados para las gráficas
del panel del tutor. Ejecuta dos consultas independientes: una sobre
el promedio ortográfico y otra sobre el promedio de legibilidad,
retornando para cada dimensión la cantidad de alumnos por categoría.

\begin{lstlisting}[language=json, caption={Estructura de respuesta del endpoint de progreso gráfico}, label={lst:progreso-response}]
/* Respuesta 200 OK */
{
  "ortografia": [
    { "categoria": "Buen Promedio",    "cantidad_alumnos": <int> },
    { "categoria": "Promedio Regular", "cantidad_alumnos": <int> },
    { "categoria": "Mal Promedio",     "cantidad_alumnos": <int> }
  ],
  "legibilidad": [
    { "categoria": "Buena Legibilidad",    "cantidad_alumnos": <int> },
    { "categoria": "Legibilidad Regular",  "cantidad_alumnos": <int> },
    { "categoria": "Mala Legibilidad",     "cantidad_alumnos": <int> }
  ]
}
\end{lstlisting}

\paragraph{Endpoints de perfil (tutor y alumno)}

\begin{table}[H]
\centering
\caption{Endpoints del módulo de perfil}
\label{tab:endpoints-perfil}
\begin{tabular}{|l|l|l|}
\hline
\textbf{Endpoint} & \textbf{Método} & \textbf{Acceso} \\
\hline
\texttt{/api/perfil/tutor/nombre}             & PUT    & Tutor (JWT) \\
\texttt{/api/perfil/tutor/correo}             & PUT    & Tutor (JWT) \\
\texttt{/api/perfil/tutor/password}           & PUT    & Tutor (JWT) \\
\texttt{/api/perfil/tutor/cuenta}             & DELETE & Tutor (JWT) \\
\texttt{/api/perfil/alumno/info}              & PUT    & Alumno (JWT) \\
\texttt{/api/perfil/alumno/password}          & PUT    & Alumno (JWT) \\
\texttt{/api/perfil/alumno/estado-tutor}      & GET    & Alumno (JWT) \\
\hline
\end{tabular}
\end{table}

El endpoint \texttt{/api/perfil/alumno/estado-tutor} merece una mención
especial: su propósito es verificar si el tutor asociado al alumno
sigue activo en el sistema (\texttt{id\_estatus = 1}). Esta verificación
protege al alumno de intentar acceder a funcionalidades dependientes
del tutor cuando este ha sido dado de baja, retornando un indicador
booleano (\texttt{tutor\_activo: true | false}) que el Front-End
utiliza para condicionar la renderización del panel del alumno.

\noindent\textbf{Nota de diseño:} todas las operaciones de actualización
de perfil que involucran datos sensibles —correo electrónico, contraseña
o eliminación de cuenta— requieren la verificación de la contraseña
actual antes de ejecutar el cambio, como capa adicional de seguridad
frente a sesiones comprometidas.